\section{问题四模型的建立与求解}
【此处插入图：问题四电价预测画像】
【此处插入图：问题四一周电价预测对照】

\subsection{建模思路}

问题四的微网物理结构继续沿用问题二、三，新的关键变化是：外部电网实时电价也成为未来未知量。每天 0:00 或日内阶段制定计划时，附件 4 中尚未发生的真实价格不能提前进入优化模型，只能用于历史训练、已发生时段更新、事后真实结算和模型评价。

因此：

\[
\boxed{
Q4\text{-}2
=
Q2
+
\text{电价预测}
+
\text{价格误差 SAA}
}
\]

以及

\[
\boxed{
Q4\text{-}3
=
Q3
+
\text{电价预测}
+
\text{日内价格更新}
+
\text{价格误差 SAA}
}.
\]

问题四在问题二、三已有负荷和光伏随机性的基础上进一步形成

\[
L,\quad PV,\quad \pi
\]

三类联合不确定量。

\subsection{电价预测与价格误差库}

设附件 4 已实现真实电价为

\[
\pi^{act}_{d,t},
\]

计划阶段中心预测为

\[
\hat\pi_{d,t},
\]

SAA 中第 \(\omega\) 个价格情景为

\[
\pi^{(\omega)}_{d,t}.
\]

必须满足

\begin{equation}
\boxed{
\text{未来 }\pi^{act}\text{ 不得提前进入计划优化}
}.
\end{equation}

电价中心预测采用与问题二预测框架一致的“同周期回溯 + 偏差校正”。每天 0:00 只使用历史真实价格预测未来 7 日：

\begin{equation}
\hat\pi_{\tau,t}^{(d)},
\qquad
\tau=d,\ldots,d+6.
\end{equation}

定义历史价格残差

\begin{equation}
e^\pi_{r,t}
=
\pi^{act}_{r,t}
-
\hat\pi_{r,t}.
\end{equation}

每个历史日保留完整 144 槽价格误差向量，并使用最近 28 个有效残差日。

问题四中的未来随机量包括

\[
L,\quad PV,\quad \pi.
\]

为尽量保留三者的历史共同结构，同一个 SAA 情景中的

\[
e^L,\quad e^{PV},\quad e^\pi
\]

尽量从同一历史日期联合抽取。

正式情景数取

\[
K=4,
\]

稳定性分析取

\[
K=8,
\]

随机种子为 2026。

价格情景为

\begin{equation}
\pi^{(\omega)}
=
\hat\pi+e^\pi.
\end{equation}

历史残差叠加后可能出现负价格，因此模型保留负价格情景，不将其人工截断为 0。

\subsection{Q4-2：考虑负荷、光伏与电价联合不确定性的 7 日滚动模型}

每天 0:00 同时预测未来负荷、光伏和电价，形成联合随机情景

\[
\left(
L^{(\omega)},
PV^{(\omega)},
\pi^{(\omega)}
\right).
\]

第一阶段变量继续为

\[
G^{plan}_{d,t},
\]

并对所有情景保持一致。

第二阶段继续沿用问题二物理执行变量

\[
G^L,\ G^{ch},\ PV^{ch},\ C,\ D,\ E,\ V,\ W,\ H,\ u.
\]

计划与执行连接为

\begin{equation}
G^L+G^{ch}+W
=
G^{plan}.
\end{equation}

对情景 \(\omega\)，定义

\[
\overline L^{(\omega)}
=
\max(L^{(\omega)}-PV^{(\omega)},0),
\]

\[
\overline{PV}^{(\omega)}
=
\max(PV^{(\omega)}-L^{(\omega)},0).
\]

物理约束为

\begin{equation}
G^L+D+H=\overline L,
\end{equation}

\begin{equation}
PV^{ch}+G^{ch}=C,
\end{equation}

\begin{equation}
PV^{ch}+V=\overline{PV},
\end{equation}

\begin{equation}
E_t
=
E_{t-1}
+
0.9C_t\Delta t
-
\frac{D_t\Delta t}{0.9},
\end{equation}

以及储电量上下限、充放电功率和充放电互斥约束。

Q4-2 的随机优化目标为

\begin{equation}
\min Z_d^{4-2}
=
\frac1K
\sum_{\omega=1}^{K}
\left[
C_d^{plan,(\omega)}
+
C_d^{em,(\omega)}
+
C_d^{future,(\omega)}
\right],
\end{equation}

其中

\begin{equation}
C_d^{plan,(\omega)}
=
\sum_t
\pi_{d,t}^{(\omega)}
G_{d,t}^{plan}
\Delta t,
\end{equation}

\begin{equation}
C_d^{em,(\omega)}
=
\sum_t
5\pi_{d,t}^{(\omega)}
H_{d,t}^{(\omega)}
\Delta t.
\end{equation}

未来 continuation cost 同样使用对应的价格情景。

最终报送费用不能使用预测价格，而应使用附件 4 的真实实现价格：

\begin{equation}
C_d^{4-2,act}
=
\sum_t
\left[
\pi_{d,t}^{act}G_{d,t}^{plan}
+
5\pi_{d,t}^{act}H_{d,t}^{act}
\right]\Delta t.
\end{equation}

\subsection{Q4-3：加入光伏预报与价格日内更新}

Q4-3 继承问题三的 0:00、6:00、12:00、18:00 四阶段结构。

0:00：

\begin{itemize}
\item 负荷采用问题二自建预测；
\item 光伏采用附件 3 的 0:00 最新预报；
\item 电价采用 0:00 价格预测。
\end{itemize}

6:00、12:00、18:00：

\begin{itemize}
\item 负荷中心预测保持当天 0:00 版本；
\item 光伏中心预测切换为当前时刻附件 3 最新预报；
\item 电价预测允许利用已经发生的真实价格进行因果更新；
\item 已执行时段保持锁定；
\item 当前真实储电量作为新阶段初值。
\end{itemize}

设当天基础价格预测为

\[
\hat\pi_{d,t}^{base}.
\]

阶段 \(s\) 只允许使用

\[
\{\pi^{act}_{d,\tau}:\tau<s\}
\]

中已经发生的真实价格构造日内校正项 \(b_{d,s}\)，则剩余时段中心预测写为

\begin{equation}
\hat\pi_{d,t}^{(d,s)}
=
\hat\pi_{d,t}^{base}
+
b_{d,s},
\qquad
t\ge s.
\end{equation}

日内更新不使用未来真实价格。

四个阶段分别维护价格残差库

\[
\mathcal E_\pi^0,\quad
\mathcal E_\pi^6,\quad
\mathcal E_\pi^{12},\quad
\mathcal E_\pi^{18},
\]

定义

\begin{equation}
e^{\pi,s}_{r,\ell}
=
\pi^{act}_{r,\ell}
-
\hat\pi^{(r,s)}_{r,\ell}.
\end{equation}

阶段 \(s\) 的联合随机情景为

\[
L^{(\omega)}
=
\hat L+e^L,
\]

\[
PV^{(\omega,s)}
=
\widehat{PV}^{(s)}
+
e^{PV,s},
\]

\begin{equation}
\pi^{(\omega,s)}
=
\hat\pi^{(s)}
+
e^{\pi,s}.
\end{equation}

\subsection{Q4-3 计划与调整变量}

0:00 原始计划：

\[
P_{d,t}.
\]

6:00、12:00、18:00 调整计划：

\[
A^6_{d,t},
\qquad
A^{12}_{d,t},
\qquad
A^{18}_{d,t}.
\]

最终生效计划为

\begin{equation}
B_{d,t}
=
\begin{cases}
P_{d,t},&0\le t<6,\\
A^6_{d,t},&6\le t<12,\\
A^{12}_{d,t},&12\le t<18,\\
A^{18}_{d,t},&18\le t<24.
\end{cases}
\end{equation}

所有调整统一相对于当天 0:00 原始计划结算：

\begin{equation}
A-P
=
\Delta^+-\Delta^-.
\end{equation}

与问题三不同，问题四的价格情景可能为负值，因此不能依靠正价格条件自动保证 \(\Delta^+\) 与 \(\Delta^-\) 单边取值，需要增加专门约束。

\subsection{H-1：负价格下的计划购电物理上界}

如果负价格情景下计划购电没有物理上界，模型可能为了获得负成本收益而购买大量没有实际用途的电量，并全部进入 \(W\)。

正常购电的实际用途只有“补负荷”和“给储能充电”，因此构造当前日计划购电物理上界

\begin{equation}
U_{d,t}
=
\max_{\omega}
\overline L_{d,t}^{(\omega)}
+
P_{\max}.
\end{equation}

即

\[
U_{d,t}
=
\max_{\omega}
\overline L_{d,t}^{(\omega)}
+
5000.
\]

于是

\[
0\le P_{d,t}\le U_{d,t},
\]

\begin{equation}
0\le A^s_{d,t}\le U^s_{d,t}.
\end{equation}

未来延续成本中的情景计划变量同样满足

\begin{equation}
0\le
G_{\tau,t}^{plan,(\omega)}
\le
\overline L_{\tau,t}^{(\omega)}
+
5000.
\end{equation}

这样 \(W\) 只表示预测误差造成的已购未用电，而不能成为负价格套利通道。

\subsection{H-2：调增与调减显式互斥}

在负价格条件下，若只设置

\[
A-P=\Delta^+-\Delta^-,
\qquad
\Delta^+,\Delta^-\ge0,
\]

可能出现两者同时为正的非业务解。

因此增加二元变量 \(z_t\)：

\[
0\le\Delta_t^+\le M_tz_t,
\]

\[
0\le\Delta_t^-\le M_t(1-z_t),
\]

\begin{equation}
z_t\in\{0,1\}.
\end{equation}

Big-\(M\) 不使用人工大数，而取

\begin{equation}
M_t=U_{d,t}.
\end{equation}

从而严格实现

\[
\Delta^+
=
\max(A-P,0),
\]

\[
\Delta^-
=
\max(P-A,0).
\]

\subsection{Q4-3 费用与目标函数}

阶段 \(s\) 的优化价格采用情景价格

\[
\pi^{(\omega,s)}.
\]

时段调整结算为

\begin{equation}
C^{settle,s,(\omega)}
=
\left[
\pi^{(\omega,s)}P
+
1.5\pi^{(\omega,s)}\Delta^+
-
0.5\pi^{(\omega,s)}\Delta^-
\right]\Delta t.
\end{equation}

紧急购电费用为

\[
5\pi^{(\omega,s)}H.
\]

0:00 目标函数为

\begin{equation}
\min Z_{d,0}^{4-3}
=
\frac1K
\sum_\omega
\left[
C^{plan,(\omega)}
+
C^{em,(\omega)}
+
C^{future,(\omega)}
\right].
\end{equation}

6:00、12:00、18:00 阶段目标函数为

\begin{equation}
\min Z_{d,s}^{4-3}
=
\frac1K
\sum_\omega
\left[
C^{settle,s,(\omega)}
+
C^{em,s,(\omega)}
+
C^{future,s,(\omega)}
\right].
\end{equation}

最终报送时全部替换为真实电价 \(\pi^{act}\)。

真实正常购电与调整费用为

\begin{equation}
C_d^{normal,act}
=
\sum_t
\left[
\pi_{d,t}^{act}P_{d,t}
+
1.5\pi_{d,t}^{act}\Delta^+_{d,t}
-
0.5\pi_{d,t}^{act}\Delta^-_{d,t}
\right]\Delta t,
\end{equation}

紧急购电费用为

\begin{equation}
C_d^{em,act}
=
\sum_t
5\pi_{d,t}^{act}
H_{d,t}^{act}\Delta t,
\end{equation}

总费用为

\begin{equation}
C_d^{4-3,act}
=
C_d^{normal,act}
+
C_d^{em,act}.
\end{equation}

\subsection{模型求解}

Q4-2 经联合情景展开后属于

\[
\boxed{\text{7 日滚动两阶段 SAA-MILP}},
\]

Q4-3 属于

\[
\boxed{\text{日内多阶段滚动 SAA-MILP}}.
\]

两者的连续能量流变量与二元充放电状态、Q4-3 的调增调减互斥状态共同构成混合整数线性规划。

底层继续采用分支定界/分支切割，并利用 MATLAB \texttt{intlinprog} 求解。首先求解 LP 松弛作为热启动，再对非整数变量进行分支，通过节点下界、整数上界和切平面逐步缩小搜索空间。

Q4-2 每天执行：

\[
\boxed{
\text{负荷预测+光伏预测+电价预测}
\rightarrow
\text{联合 SAA 情景}
\rightarrow
\text{7 日滚动 MILP}
\rightarrow
\text{输出当天计划}
\rightarrow
\text{真实执行并更新储电量}
}
\]

Q4-3 日内执行：

\[
\boxed{
\begin{array}{c}
0{:}00:\ \text{负荷预测+光伏预报+价格预测，求 }P\\
\downarrow\\
\text{执行至 }6{:}00\\
\downarrow\\
\text{读取新光伏预报+已实现价格+真实储电量，求 }A^6\\
\downarrow\\
\text{执行至 }12{:}00,\ \text{求 }A^{12}\\
\downarrow\\
\text{执行至 }18{:}00,\ \text{求 }A^{18}\\
\downarrow\\
\text{执行至 }24{:}00
\end{array}
}
\]

即坚持

\[
\boxed{
\text{求解}
\rightarrow
\text{执行}
\rightarrow
\text{更新真实状态和价格信息}
\rightarrow
\text{再求解}
}.
\]

\subsection{结果分析框架}



\begin{table}[H]
\centering
\zihao{5}
\begin{tabular}{lrrrr}
\toprule
指标 & Q2 固定电价 & Q4-2 波动电价 & Q3 固定电价 & Q4-3 波动电价 \\
\midrule
总费用/元 & 14806890.28 & 16320514.83 & 14485230.30 & 15225713.07 \\
正常+调整费用/元 & — & 13737348.89 & 13451415.12 & 14148292.42 \\
紧急购电费用/元 & 1657825.30 & 2583165.94 & 1033815.18 & 1077420.65 \\
紧急购电量/kWh & 269240.94 & 450371.13 & 185523.91 & 200511.20 \\
已购未用电/kWh & 1136190.28 & 1061902.30 & 601472.99 & 566918.30 \\
日末储电量均值/kWh & 7045.97 & 4810.16 & 6847.95 & 4831.70 \\
\bottomrule
\end{tabular}
\caption{问题四两口径费用与紧急购电对照（费用单位：元，电量单位：kWh）}
\label{tab:11}
\end{table}

结果分析应重点回答：

\begin{enumerate}
\item 波动价格相对于固定价格如何改变总成本；
\item 储能是否表现出更明显的低价充电、高价放电；
\item 电价预测误差造成多少额外成本；
\item Q4-3 的日内光伏与价格更新是否进一步降低紧急购电或调整成本；
\item 负价格情景下 H-1/H-2 是否被实际激活；
\item 价格峰谷预测偏差如何影响储能调度时机。
\end{enumerate}

\subsection{电价预测与价格风险检验}

电价预测精度使用

\[
MAE_\pi
=
\frac1N
\sum_i
|\hat\pi_i-\pi_i^{act}|,
\]

\begin{equation}
RMSE_\pi
=
\sqrt{
\frac1N
\sum_i
(\hat\pi_i-\pi_i^{act})^2
}.
\end{equation}

\begin{table}[H]
\centering
\zihao{5}
\begin{tabular}{lr}
\toprule
指标 & 数值 \\
\midrule
MAE & 0.0432 \\
RMSE & 0.0602 \\
峰值时刻精确命中率 & 36.2\% \\
谷值时刻精确命中率 & 33.8\% \\
峰值 1 h 内命中率 & 63.5\% \\
谷值 1 h 内命中率 & 58.4\% \\
峰谷差预测 MAE & 0.0675 \\
\bottomrule
\end{tabular}
\caption{电价预测精度指标（误差单位：元/kWh）}
\label{tab:12}
\end{table}

同时比较只使用中心价格预测的方案 \(P0\) 与“中心预测 + 价格误差 SAA”的方案 \(P1\)，衡量价格风险建模的贡献。

建立“未来真实价格提前已知”的完美价格信息基准，用于计算价格不确定性代价：

\begin{equation}
\Delta C_{\rm price}
=
C_{\rm forecast}
-
C_{\rm perfect-price}.
\end{equation}

其结果为

\[
\Delta C_{\rm price}
=
\text{66993.37 元}.
\]

\subsection{模型检验}

问题四除继续检查能量平衡、储能递推、充放电互斥、正常购电守恒和非预见性外，还应检查：

\begin{itemize}
\item 未来真实价格无泄漏；
\item 三类误差同日联合抽样；
\item Q4-3 阶段价格更新的因果性；
\item 最终费用按真实价格重算；
\item 紧急购电使用 \(5\pi^{act}\)；
\item 调增/调减使用 \(1.5\pi^{act}\) 与 \(0.5\pi^{act}\)；
\item H-1 计划购电物理上界；
\item H-2 调增/调减互斥；
\item 负价格情景不导致模型无界；
\item \(W\) 不出现负价套利；
\item 阶段执行与最终重放一致。
\end{itemize}

H-2 启用后，应有

\[
\#\{t:\Delta_t^+>0\text{ 且 }\Delta_t^->0\}=0.
\]

最大等式残差：

\[
6.96\times10^{-10}.
\]

最大整数间隙：

\[
1.16\times10^{-3}.
\]

最大拼接重放偏差：

\[
0.
\]

稳定性分析比较 \(K=4\) 与 \(K=8\)：

\begin{table}[H]
\centering
\zihao{5}
\begin{tabular}{lrrr}
\toprule
指标 & \(K=4\) & \(K=8\) & 相对差异 \\
\midrule
Q4-2 总费用 & 16320514.83 & 16240197.29 & $-0.49\%$ \\
Q4-3 总费用 & 15225713.07 & 15100131.61 & $-0.82\%$ \\
Q4-2 紧急购电量 & 450371.13 & 441055.51 & $-2.07\%$ \\
Q4-3 紧急购电量 & 200511.20 & 186767.09 & $-6.85\%$ \\
\bottomrule
\end{tabular}
\caption{问题四 Q4-2 与 Q4-3 指标对照}
\label{tab:13}
\end{table}
